% ============================================================
%  Chapter 4: Experimental Setup
%  File: chapters/04_experiments.tex
%  Bibliography: chapters/refs_04_setup.bib
% ============================================================

\chapter{Experimental Setup}
\label{ch:experiments}

This chapter describes the experimental protocol used to evaluate
\modnet{}. We begin with the sixteen benchmark problems, then describe
the data encoding, the evaluation procedure, the hyperparameter choices,
the visualization method, and finally the reproducibility of the
experiments.

% ------------------------------------------------------------
\section{Benchmark Problems}
\label{sec:exp-tasks}
% ------------------------------------------------------------

We evaluate \modnet{} on sixteen tasks that span four broad classes:

\begin{enumerate}
  \item \textbf{Boolean logic:} XOR-2, XOR-3, Parity-4, Parity-6.
  \item \textbf{Majority functions:} Majority-3, Majority-5.
  \item \textbf{Digital arithmetic and circuits:} Full Adder, 2-bit
    Adder, Comparator, 4-to-1 Multiplexer, 4-bit Sorting Network.
  \item \textbf{Continuous regression:} $\sin(2\pi x)$, $\sin(4\pi x)$,
    $x^2 + y^2$, Gaussian, Multiplication.
\end{enumerate}

The tasks were chosen to cover a spectrum of difficulty, from problems
that are trivially solvable by a single linear unit (XOR-2) to problems
that require multi-stage nonlinear computation and non-monotone Boolean
logic (Parity-6, 4-bit sorting). Table~\ref{tab:all-tasks} summarizes
the tasks.

\begin{table}[h]
  \centering
  \small
  \caption{Summary of the sixteen benchmark problems.}
  \label{tab:all-tasks}
  \begin{tabular}{llrrrl}
    \toprule
    \textbf{Task} & \textbf{Class} & $n_{\text{in}}$ & $n_{\text{out}}$
      & $N$ & \textbf{Description} \\
    \midrule
    XOR-2        & Boolean    & 2 & 1 &  4 & $a \oplus b$ \\
    XOR-3        & Boolean    & 3 & 1 &  8 & $a \oplus b \oplus c$ \\
    Parity-4     & Boolean    & 4 & 1 & 16 & $\bigoplus_{i=1}^4 x_i$ \\
    Parity-6     & Boolean    & 6 & 1 & 64 & $\bigoplus_{i=1}^6 x_i$ \\
    \midrule
    Majority-3   & Boolean    & 3 & 1 &  8 & $\sum x_i \geq 2$ \\
    Majority-5   & Boolean    & 5 & 1 & 32 & $\sum x_i \geq 3$ \\
    \midrule
    Full Adder   & Boolean    & 3 & 2 &  8 & $s = a \oplus b \oplus c$, $c_{\text{out}} = a b + b c + a c$ \\
    2-bit Adder  & Boolean    & 4 & 3 & 16 & $(a_0 + 2 a_1) + (b_0 + 2 b_1)$ \\
    Comparator   & Boolean    & 2 & 3 &  4 & $(\text{gt}, \text{eq}, \text{lt})$ one-hot \\
    MUX 4-to-1   & Boolean    & 6 & 1 & 64 & $d_{s_0 + 2 s_1}$ \\
    Sort-4       & Boolean    & 4 & 4 & 16 & $(\min, \ldots, \max)$ of 4 bits \\
    \midrule
    $\sin(2\pi x)$     & Regression & 4 & 8 & 32 & $0.5 + 0.5 \sin(2\pi x)$ \\
    $\sin(4\pi x)$     & Regression & 4 & 8 & 32 & $0.5 + 0.5 \sin(4\pi x)$ \\
    $x^2 + y^2$        & Regression & 8 & 8 & 25 & $0.5 + 0.25 (x_1^2 + x_2^2)$ \\
    Gaussian           & Regression & 4 & 8 & 25 & $\exp(-x^2)$ on $[-1,1]$ \\
    Multiplication     & Regression & 8 & 8 & 36 & $x \cdot y$ for $x,y \in [0,1]$ \\
    \bottomrule
  \end{tabular}
  \vspace{0.3em}
  \par\footnotesize{$N$ denotes the number of training samples.}
\end{table}

\subsection{Boolean Function Learning}
\label{sec:exp-boolean}

Boolean function learning is a natural testbed because the fitness
function is exact (no measurement noise) and correctness is unambiguous.

\paragraph{XOR-2 and XOR-3.} The exclusive-or of two or three bits. XOR
is the canonical example of a non-linearly-separable Boolean function:
no single linear threshold unit can compute it
\citep{minsky1969perceptrons}.

\paragraph{Parity-4 and Parity-6.} The parity of $n$ bits,
$\text{Parity}_n(\mathbf{x}) = \bigoplus_{i=1}^{n} x_i$. Parity is the
XOR of $n$ bits and is the simplest Boolean function that is provably
hard for bounded-depth threshold circuits: any depth-$d$ circuit that
computes parity requires $\Omega(n^{1+\varepsilon_d})$ wires
\citep{impagliazzo1997size}. The number of input--output pairs grows
exponentially with $n$.

\paragraph{Majority-3 and Majority-5.} The majority function returns
$1$ if and only if the number of active bits exceeds $n/2$. Unlike
parity, majority is monotone: flipping any input from $0$ to $1$
cannot decrease the output.

\subsection{Digital Circuits}
\label{sec:exp-circuits}

The digital circuit tasks exercise \modnet{} on problems that are
canonical in hardware design. Each is a well-defined Boolean function
with a known minimal circuit.

\paragraph{Full Adder.} Three inputs $(a, b, c_{\text{in}})$, two
outputs $(\text{sum}, c_{\text{out}})$. The sum is XOR of the three
inputs; the carry-out is the majority.

\paragraph{2-bit Adder.} Four inputs $(a_0, a_1, b_0, b_1)$, three
outputs $(s_0, s_1, c_{\text{out}})$. Represents the addition of two
2-bit numbers.

\paragraph{Comparator.} Two inputs $(a, b)$, three one-hot outputs:
whether $a > b$, $a = b$, or $a < b$.

\paragraph{4-to-1 Multiplexer.} Six inputs (four data lines
$d_0, \ldots, d_3$, two select lines $s_0, s_1$), one output: the data
line selected by the two select bits.

\paragraph{4-bit Sorting Network.} Four inputs, four outputs: the
inputs sorted in ascending order. This is nontrivial because sorting
is \emph{non-monotone} in each input bit.

\subsection{Continuous Regression}
\label{sec:exp-regression}

Regression tasks test whether \modnet{} can produce smooth, graded
outputs rather than binary decisions.

\paragraph{$\sin(2\pi x)$ and $\sin(4\pi x)$.} Two sinusoids over the
unit interval, normalized to $[0, 1]$. The second has twice the
frequency of the first.

\paragraph{$x^2 + y^2$.} A two-dimensional quadratic surface. We sample
a $5 \times 5$ grid over $[0,1]^2$.

\paragraph{Gaussian.} The function $\exp(-x^2)$ on $[-1, 1]$, normalized
to $[0, 1]$.

\paragraph{Multiplication.} The product $x \cdot y$ for $x, y \in [0,1]$.
We sample a $6 \times 6$ grid, giving $36$ points.

% ------------------------------------------------------------
\section{Data Encoding}
\label{sec:exp-encoding}
% ------------------------------------------------------------

Inputs and outputs of \modnet{} are binary vectors. Continuous values
are quantized to a fixed number of bits.

\subsection{Boolean Inputs}

For Boolean tasks, each input bit is passed directly to the network as
a $0$ or $1$ value. No encoding is required.

\subsection{Continuous Inputs}

For regression tasks, each continuous input variable $x \in [0, 1]$
is quantized to $b$ bits as follows. Let
$v = \lfloor x \cdot (2^b - 1) \rceil$ be the rounded integer in
$[0, 2^b - 1]$. The bit vector is then

\begin{equation}
  \text{encode}(x, b) =
    \left(
      \lfloor v / 2^0 \rfloor \bmod 2,\,
      \lfloor v / 2^1 \rfloor \bmod 2,\,
      \ldots,\,
      \lfloor v / 2^{b-1} \rfloor \bmod 2
    \right).
  \label{eq:encode}
\end{equation}

We use $b = 4$ bits for inputs, giving $2^4 = 16$ distinct input
levels per variable. This is a coarse quantization, deliberately
chosen to keep the input space small enough for exhaustive evaluation.

\subsection{Outputs}

For Boolean tasks, outputs are single bits or short bit-vectors
(for multi-output circuits, up to $4$ bits). For regression tasks,
outputs are quantized to $b = 8$ bits, giving $2^8 = 256$ distinct
output levels. The decoding function is the inverse of
Equation~\ref{eq:encode}:

\begin{equation}
  \text{decode}(\mathbf{y}) =
    \frac{1}{2^b - 1} \sum_{i=0}^{b-1} y_i \cdot 2^i.
  \label{eq:decode}
\end{equation}

The choice of $8$ bits for outputs is a trade-off between resolution
and computational cost. With $8$ bits, the quantization error is at
most $1/255 \approx 0.004$, well below the RMSE values we report.

% ------------------------------------------------------------
\section{Evaluation Procedure}
\label{sec:exp-eval}
% ------------------------------------------------------------

Each network is evaluated on the full training set
$\mathcal{D} = \{(\mathbf{x}^{(k)}, \mathbf{y}^{(k)})\}_{k=1}^{K}$.
The network is initialized with all node activations set to zero, the
input is applied at the designated input nodes, and the network is run
for $T$ time steps. The output is read from the output nodes.

The fitness value is computed using the mean squared error between
the predicted output and the target. For single-output Boolean tasks
and all regression tasks, this is a scalar comparison; for
multi-output Boolean tasks, it is a bit-level comparison.

An experiment is considered \emph{solved} when:

\begin{itemize}
  \item \textbf{Boolean tasks.} All $K$ training samples are classified
    correctly (exact $100\%$ accuracy).
  \item \textbf{Regression tasks.} The root-mean-square error is below
    a task-specific threshold. We use RMSE $< 0.05$ for all regression
    tasks except $\sin(4\pi x)$, for which the higher frequency
    requires RMSE $< 0.10$.
\end{itemize}

% ------------------------------------------------------------
\section{Hyperparameter Settings}
\label{sec:exp-hp}
% ------------------------------------------------------------

The hyperparameters of \modnet{} are listed in
Table~\ref{tab:hp-setup}. Unless otherwise noted, all tasks use the
same values.

\begin{table}[h]
  \centering
  \small
  \caption{Hyperparameters of \modnet{} used in all experiments.}
  \label{tab:hp-setup}
  \begin{tabular}{lll}
    \toprule
    \textbf{Parameter} & \textbf{Description} & \textbf{Value} \\
    \midrule
    $P$                & Population size                    & 300 \\
    $G$                & Maximum generations                & 400--1500 \\
    $p_{\text{cross}}$ & Crossover rate                     & 0.30 \\
    $p_{\text{fresh}}$ & Fresh random rate                  & 0.05 \\
    $\sigma_{\text{base}}$ & Base mutation scale           & 0.30 \\
    $\sigma_{\max}$    & Maximum mutation scale             & 1.50 \\
    $\gamma$           & Mutation growth factor             & 1.25 \\
    $g_{\text{thresh}}$& Plateau threshold (generations)    & 25 \\
    $g_{\text{birth}}$ & Birth interval (generations)       & 100 \\
    $g_{\text{prune}}$ & Prune interval (generations)       & 20 \\
    $\epsilon_{\text{prune}}$ & Prune tolerance             & $-0.001$ \\
    $w_{\text{thresh}}$& Weight pruning threshold           & 0.04 \\
    $\alpha_{\text{lo}}$& Lower activity threshold          & 0.02 \\
    $\alpha_{\text{hi}}$& Upper activity threshold          & 0.98 \\
    $n_{\text{hidden,min}}$& Minimum hidden nodes          & 2 \\
    $n_{\max}$         & Maximum nodes                      & 60 \\
    \bottomrule
  \end{tabular}
\end{table}

Task-specific settings of the number of time steps $T$, the number of
initial nodes, and the maximum number of generations are listed in
Table~\ref{tab:task-hp}.

\begin{table}[h]
  \centering
  \small
  \caption{Task-specific hyperparameters. $T$ is the number of time
    steps per evaluation, $n_0$ is the initial number of nodes,
    $G_{\max}$ is the maximum number of generations.}
  \label{tab:task-hp}
  \begin{tabular}{lrrr}
    \toprule
    \textbf{Task} & $T$ & $n_0$ & $G_{\max}$ \\
    \midrule
    XOR-2         & 4 & 26 &  400 \\
    XOR-3         & 5 & 26 &  600 \\
    Parity-4      & 6 & 26 &  800 \\
    Parity-6      & 8 & 26 & 1500 \\
    Majority-3    & 5 & 26 &  500 \\
    Majority-5    & 7 & 26 &  900 \\
    Full Adder    & 6 & 26 &  800 \\
    2-bit Adder   & 8 & 26 & 1200 \\
    Comparator    & 6 & 26 &  600 \\
    MUX 4-to-1    & 8 & 26 & 1200 \\
    Sort-4        & 8 & 26 & 1200 \\
    $\sin(2\pi x)$& 5 & 26 & 1500 \\
    $\sin(4\pi x)$& 6 & 26 & 1500 \\
    $x^2 + y^2$   & 6 & 26 & 1200 \\
    Gaussian      & 5 & 26 & 1500 \\
    Multiplication& 6 & 26 & 1500 \\
    \bottomrule
  \end{tabular}
\end{table}

The number of time steps $T$ is set based on the depth of the required
computation. Boolean tasks with more inputs require more steps to
propagate information through the network; regression tasks with
periodic outputs require additional steps to resolve the periodicity.

All experiments start from an initial network of $n_0 = 26$ nodes
(one output for single-output tasks, eight outputs for regression, two
to four outputs for multi-output Boolean). All nodes may connect to
any other node, including themselves.

% ------------------------------------------------------------
\section{Network Visualization}
\label{sec:exp-viz}
% ------------------------------------------------------------

To make the evolved topologies interpretable, we render every network
using a \emph{force-directed} (spring) layout implemented in the
\texttt{networkx} library \citep{hagberg2008exploring}. The spring
layout treats each edge as an attractive force and each node as a
repulsive force, causing the graph to settle into a stable
configuration in which topologically related nodes appear nearby.

Nodes are colored by role:

\begin{itemize}
  \item \textbf{Green squares:} input ports, labeled $x_0, x_1, \ldots$
  \item \textbf{Blue circles:} hidden nodes, labeled $n_i$
  \item \textbf{Red circles:} output nodes, labeled $n_i$ and marked
    with an ``out'' arrow
\end{itemize}

Edges are colored by type:

\begin{itemize}
  \item \textbf{Green:} input edges (from input ports to nodes)
  \item \textbf{Blue:} internal edges (between nodes)
  \item \textbf{Red arcs:} self-loops (recurrent connections)
\end{itemize}

This visualization is used throughout Chapter~\ref{ch:results} to
present the evolved networks. Figure~\ref{fig:viz-example} shows an
example for the XOR-2 task.

\begin{figure}[h]
  \centering
  \safeincludegraphics[width=0.6\textwidth]{xor2/best_spring.png}
  \caption{Spring-layout visualization of the XOR-2 network. Green
    squares are input ports, blue circles are hidden nodes, red circles
    are output nodes. The single self-loop is shown as a red arc. The
    network has $8$ nodes and $17$ active connections.}
  \label{fig:viz-example}
\end{figure}

% ------------------------------------------------------------
\section{Reproducibility}
\label{sec:exp-repro}
% ------------------------------------------------------------

All experiments use a fixed random seed of $42$, reset before each
task to ensure that the results for one task do not depend on the
order of execution. The seed is used to initialize the population and
to control all subsequent stochastic operations (mutation, crossover,
and selection of parents).

Each experiment produces the following outputs:

\begin{itemize}
  \item \texttt{best.json}: the complete evolved network, including
    all weights, biases, and node roles.
  \item \texttt{best\_spring.png}: a spring-layout rendering of the
    final network.
  \item \texttt{snapshot\_gen\_XXXX.png}: snapshots of the best
    network at selected generations, showing how the structure
    evolved.
  \item \texttt{curve.csv} and \texttt{curve.svg}: the fitness history
    over all generations (where applicable).
  \item \texttt{samples.csv}: the network's predictions on all
    training samples.
  \item \texttt{events.log} and \texttt{events.csv}: a complete log
    of every birth and prune operation, with timestamps (where
    applicable).
\end{itemize}

The complete experimental framework is implemented in a single Python
script and is available as open-source software
\citep{modnet2026code}. The script has no external dependencies beyond
NumPy, and optionally \texttt{networkx} and \texttt{matplotlib} for
visualization. It runs on commodity hardware, including mobile phones.

% ------------------------------------------------------------
\section{Computational Cost}
\label{sec:exp-cost}
% ------------------------------------------------------------

The total wall-clock time for all sixteen experiments is
approximately $320$ seconds on a single phone CPU (ARM Cortex-A78,
$2.4$~GHz, $8$ cores). The breakdown is shown in
Table~\ref{tab:time}.

\begin{table}[h]
  \centering
  \small
  \caption{Wall-clock time per task, sorted by ascending time.}
  \label{tab:time}
  \begin{tabular}{lrr}
    \toprule
    \textbf{Task} & \textbf{Time (s)} & \textbf{Generations} \\
    \midrule
    XOR-2         &  0.8 &   43 \\
    Majority-3    &  0.9 &   44 \\
    Parity-4      &  1.8 &   99 \\
    XOR-3         &  2.0 &  130 \\
    Full Adder    &  3.0 &  156 \\
    Comparator    &  3.8 &  251 \\
    MUX 4-to-1    &  4.8 &  118 \\
    2-bit Adder   &  7.3 &  350 \\
    Sort-4        &  8.7 &  417 \\
    Majority-5    & 10.2 &  291 \\
    $x^2 + y^2$   & 21.1 &  862 \\
    $\sin(4\pi x)$& 25.1 & 1499 \\
    Multiplication& 28.5 & 1499 \\
    $\sin(2\pi x)$& 30.0 & 1499 \\
    Parity-6      & 40.6 & 1499 \\
    Gaussian      & 41.6 & 1499 \\
    \midrule
    \textbf{Total} & \textbf{320.2} & \\
    \bottomrule
  \end{tabular}
\end{table}

The time scales roughly with the difficulty of the task: XOR-2 and
Majority-3 finish in under one second, while Parity-6, Gaussian, and
the sinusoid tasks require tens of seconds. The longest-running
tasks use the highest number of generations and the largest evolved
networks.

The computational cost of a single evaluation is $O(E \cdot T)$, where
$E$ is the number of active connections and $T$ is the number of time
steps. For the largest network in our experiments ($N = 27$ nodes,
$E \approx 350$ active connections, $T = 5$ time steps), a single
evaluation takes approximately $50\,\mu$s. Since the evolutionary
algorithm performs $P \cdot G \approx 300 \cdot 1500 = 4.5 \times
10^5$ evaluations, the total compute for the longest task is on the
order of $20$ seconds of pure evaluation time.

% ------------------------------------------------------------
\section{Summary of Reproducibility Artifacts}
\label{sec:exp-artifacts}
% ------------------------------------------------------------

Table~\ref{tab:artifacts} lists the reproducibility artifacts
produced by each experiment.

\begin{table}[h]
  \centering
  \small
  \caption{Reproducibility artifacts per task.}
  \label{tab:artifacts}
  \begin{tabular}{lp{9cm}}
    \toprule
    \textbf{Artifact} & \textbf{Description} \\
    \midrule
    \texttt{best.json}         & Complete network state (weights,
      biases, output indices) in JSON format. \\
    \texttt{best\_spring.png}  & Force-directed (spring) layout of
      the final network. \\
    \texttt{snapshot\_gen\_XXXX.png} & Snapshots of the best network
      at five selected generations. \\
    \texttt{curve.csv}         & Fitness, recursive count, node count,
      and wire count at each generation. \\
    \texttt{samples.csv}       & Predictions on all training samples. \\
    \texttt{events.log}        & Human-readable log of every birth and
      prune operation. \\
    \bottomrule
  \end{tabular}
\end{table}

These artifacts allow full reproduction and inspection of every
evolved network, and were used to generate all figures and tables
reported in Chapter~\ref{ch:results}.

% ============================================================
%  Bibliography for this chapter
% ============================================================
